% GENERATED FILE. Edit canonical lessons, then run npm run generate:books.
% canonical-source-hash: fnv1a64:a8e95859f1ba4235
% canonical-lessons: FR-C264-jetais-content, FR-C264-jespere, FR-C264-revision-sentiments

\chapter{How It Felt — j'étais content, j'espère}
\label{ch:fr-how-it-felt}
\hlchaptermodality{264}

\begin{chapteropening}
\textbf{By the end of this chapter:} \emph{I can say how an experience felt, with the feeling as scene, and say what I hope and dream of doing next.}

Complete FR-C264-revision-sentiments: how a journey felt, and what you hope to do next.
\end{chapteropening}

\section[j'étais content, j'étais inquiet]{j'étais content — how it felt}

\label{lesson:FR-C264-jetais-content}



\begin{quote}
Say the last line of the journey you read: \textquotedblleft{}I was pleased\textquotedblright{}. (\emph{J'étais content.})

That sentence chose its past for a reason. This lesson is about why.
\end{quote}

\begin{grammarlens}[title={Grammar Lens: a feeling is scene}]
You already own the feelings: \textbf{content}, \textbf{heureux}, \textbf{triste}, \textbf{inquiet}, \textbf{nerveux}, \textbf{fier}. A feeling has no edge. It was going on, like the weather, so in a story it is scene, and it takes the imparfait of \emph{être}:

\begin{quote}
\emph{D'abord, j'étais nerveux. Ensuite, j'étais content.} — At first I was nervous. Then I was pleased.
\end{quote}

The adjective agrees with whoever felt it, and the ear hears some of that. \emph{Content} ends in a silent \emph{t}; \textbf{contente} sounds it. \emph{Inquiet} becomes \textbf{inquiète}, \emph{nerveux} becomes \textbf{nerveuse}, \emph{heureux} becomes \textbf{heureuse}. \emph{Triste} has one form for a man or a woman.

\begin{quote}
\emph{Ma mère était inquiète.} — My mother was worried. \emph{Nous étions contents.} — We were pleased.
\end{quote}

The passé composé is possible too, and it changes the picture: it makes the feeling a reaction, an event at one moment. \emph{Quand j'ai vu mon ami, j'ai été content.} — When I saw my friend, I was glad.
\end{grammarlens}

\begin{grammarlens}[title={Grammar Lens: content de}]
To say \textbf{what} you were pleased with, add \textbf{de}:

\begin{quote}
\emph{J'étais content \textbf{du} voyage.} — I was pleased with the trip. \emph{Je suis content \textbf{de} partir.} — I'm glad to be leaving.
\end{quote}

\emph{De} and \emph{le} weld into \textbf{du}, and an infinitive can follow \emph{de} directly. \textbf{Fier} works the same way: \emph{je suis fier de toi}, I'm proud of you.
\end{grammarlens}

\begin{cousinweb}
\textbf{content} is usually traced to Latin \emph{contentus}, \textquotedblleft{}satisfied\textquotedblright{}, the participle of \emph{continēre}, \textquotedblleft{}to hold together\textquotedblright{} (Lewis \& Short, under \emph{contentus} and \emph{contineo}), itself \emph{con-} plus \emph{tenēre}, \textquotedblleft{}to hold\textquotedblright{}. EtymDB 2.1 follows it back through Old French \emph{content}. To be content is to hold yourself in, wanting nothing further. \textbf{inquiet} is a borrowing from Latin \emph{inquiētus}, \textquotedblleft{}restless\textquotedblright{} (EtymDB 2.1; Lewis \& Short): \emph{in-}, \textquotedblleft{}not\textquotedblright{}, and \emph{quiētus}, \textquotedblleft{}quiet\textquotedblright{}.
\end{cousinweb}

\subsection*{Your turn}
Say these aloud:

\begin{itemize}
\raggedright
  \item \emph{D'abord, j'étais nerveux. Ensuite, j'étais content.}
  \item the pairs — \emph{content, contente} — \emph{inquiet, inquiète}
  \item \emph{J'étais content du voyage.}
  \item a feeling of yours, last week, with its reason — \emph{J'étais … parce que …}
\end{itemize}

\begin{tcolorbox}[breakable,colback=teal!4,colframe=teal!35!black,title={Before you move on}]
Why is \emph{j'étais content} the usual choice in a story? (A feeling has \textbf{no edge}: it is scene.) What does \emph{j'ai été content} make of it? (A \textbf{reaction}, an event.) Say \textquotedblleft{}I was pleased with the trip\textquotedblright{}. (\emph{\textbf{J'étais content du voyage.}}) Which Latin verb hides in \emph{content}? (\emph{\textbf{Tenēre}}, \textquotedblleft{}to hold\textquotedblright{}.)
\end{tcolorbox}
\section[j'espère revenir]{j'espère revenir — saying what you are reaching for}

\label{lesson:FR-C264-jespere}



\begin{quote}
Say \textquotedblleft{}to hope\textquotedblright{}. (\emph{Espérer.}) Then \textquotedblleft{}I was pleased\textquotedblright{}. (\emph{J'étais content.})

An experience has a past and a future. Last lesson was how it felt; this one is what you want next.
\end{quote}

\subsection*{You'll want to know: j'espère + an infinitive}
\textbf{j'espère} + infinitive — \textquotedblleft{}I hope to\textquotedblright{}.

\begin{quote}
\textbf{J'espère revenir} à Paris. — I hope to come back to Paris. \textbf{J'espère parler} français un jour. — I hope to speak French one day.
\end{quote}

\textbf{espérer} is an \emph{-er} verb whose \emph{é} opens to \emph{è} where the ending goes quiet: \emph{j'espère}, \emph{tu espères}, but \emph{nous espérons}. An infinitive straight after it, with nothing in between, makes the hope your own. Italian puts a small \emph{di} there; French puts nothing.

When someone else is the one who will do it, the frame is \textbf{j'espère que…}, and the verb after it stays plain: \emph{J'espère que tu vas bien.} — I hope you're well.

\begin{grammarlens}[title={Grammar Lens: mon rêve, c'est de…}]
An ambition is bigger than a hope, and it has its own frame:

\begin{quote}
\textbf{Mon rêve, c'est de} vivre près de la mer. — My dream is to live by the sea.
\end{quote}

You have \textbf{le rêve}, \textquotedblleft{}the dream\textquotedblright{}, and \textbf{rêver}, \textquotedblleft{}to dream\textquotedblright{}, which takes the same \emph{de}: \emph{je rêve de partir}, I dream of leaving.

Three strengths, then, for the future: \emph{\textbf{je voudrais} partir}, I'd like to; \emph{\textbf{j'espère} partir}, I hope to; \emph{\textbf{mon rêve, c'est de} partir}, it is my dream. Only the last needs a \emph{de}.
\end{grammarlens}

\begin{cousinweb}
\textbf{espérer} is inherited from Latin \emph{spērāre}, \textquotedblleft{}to hope, to expect\textquotedblright{} (EtymDB 2.1; Lewis \& Short, under \emph{spero}), with an \emph{e-} that French grew in front of \emph{sp-}, as it grew an \emph{é-} in front of the \emph{sc-} of \emph{écrire}. English \textbf{desperate} is the same verb with \emph{dē-} in front: hope taken away. \textbf{rêver} goes back to Old French \emph{resver}, \textquotedblleft{}to be delirious\textquotedblright{} (EtymDB 2.1), and further than that the sources do not reach.
\end{cousinweb}

\subsection*{Your turn}
Say these aloud:

\begin{itemize}
\raggedright
  \item \emph{J'espère revenir à Paris.}
  \item someone else — \emph{J'espère que tu vas bien.}
  \item \emph{Mon rêve, c'est de vivre près de la mer.}
  \item the three strengths — \emph{je voudrais}, \emph{j'espère}, \emph{mon rêve, c'est de} — with one verb of your own
\end{itemize}

\begin{tcolorbox}[breakable,colback=teal!4,colframe=teal!35!black,title={Before you move on}]
What goes between \emph{j'espère} and the infinitive? (\textbf{Nothing.}) Which of the three strengths needs a \emph{de}? (\emph{\textbf{Mon rêve, c'est de.}}) Which English word is \emph{spērāre} with \emph{dē-} in front? (\textbf{Desperate.})
\end{tcolorbox}
\section[(review)]{Review — how it felt, and what comes next}

\label{lesson:FR-C264-revision-sentiments}



\begin{quote}
Say \textquotedblleft{}my mother was worried\textquotedblright{}. (\emph{Ma mère était inquiète.}) Then \textquotedblleft{}I hope to come back\textquotedblright{}. (\emph{J'espère revenir.})
\end{quote}

\begin{grammarlens}[title={What you've built}]
Two lessons, and an experience has a shape: what happened, how it felt, and what you want next. The feeling is scene, in the imparfait: \emph{j'étais content}, \emph{j'étais nerveux}. What it was about takes \emph{de}: \emph{content du voyage}. And the future comes in three strengths: \emph{je voudrais}, \emph{j'espère}, \emph{mon rêve, c'est de}.

Here is a journey told with all three:

\begin{quote}
\emph{D'abord, j'étais nerveux, parce que je ne connaissais pas la ville. Mais ensuite, j'étais content du voyage. J'espère revenir la semaine prochaine.}
\end{quote}

\emph{Je ne connaissais pas} is \emph{connaître} in the imparfait, on the \emph{nous connaissons} stem: a not-knowing that went on, so it is scene.
\end{grammarlens}

\subsection*{Your turn}
Say these aloud:

\begin{itemize}
\raggedright
  \item \emph{D'abord, j'étais nerveux.}
  \item \emph{Mais ensuite, j'étais content du voyage.}
  \item \emph{J'espère revenir la semaine prochaine.}
  \item a feeling with an arrow — \emph{J'étais fatigué, donc je suis rentré.}
  \item one journey of yours — how it felt at first, how it felt at the end, and one hope
\end{itemize}

\begin{tcolorbox}[breakable,colback=teal!4,colframe=teal!35!black,title={Before you move on}]
Is \emph{j'étais nerveux} scene or event? (\textbf{Scene.}) What goes between \emph{j'espère} and \emph{revenir}? (\textbf{Nothing.}) Say \textquotedblleft{}my dream is to live by the sea\textquotedblright{}. (\emph{\textbf{Mon rêve, c'est de vivre près de la mer.}})
\end{tcolorbox}
